%============================================================
% CH2 · P5 · PROBLEM 2 — SOLVE THE MAZE (Codeforces 1365D)
%============================================================
\subsection{Problem 2.5.2: Solve The Maze}
\problemlabel{2.5.2}{ssec:ch2-solve-maze}

% cells with a text: green G reached / not reached, visit numbers
\tikzset{
    g2p5-Gn/.style={g2p5-e, fill=basefill, draw=successgreen, line width=1.1pt},
    g2p5-Bn/.style={g2p5-e, fill=dupfill, draw=recursionred, line width=1.1pt},
    g2p5-on/.style={g2p5-e, fill=computedfill}
}
% \gptworow{y}{style/text, ...}: one row of cells at height y
\newcommand{\gptworow}[2]{\foreach \s/\t [count=\c from 0] in {#2} \node[g2p5-\s] at (\c,-#1) {\t};}

\problemcard{2.5.2}{Solve The Maze}{Codeforces 1365D}{https://codeforces.com/problemset/problem/1365/D}{\exhard\ \cf{1700}}{Wall in every bad person, then one search from the exit must reach every good person and no bad one}

\subsubsection*{1. Problem statement}

A maze is an $n \times m$ grid of cells: empty (\texttt{.}), wall (\texttt{\#}), good person (\texttt{G}) or bad person (\texttt{B}). People move between side-adjacent non-wall cells; the only exit is the cell $(n, m)$, which is guaranteed to be empty. You may turn any set of \emph{empty} cells (the exit included) into walls; \texttt{G} and \texttt{B} cells cannot be blocked and can be walked through. Can you do it so that every good person can reach the exit and no bad person can? Print \texttt{Yes} or \texttt{No}.

\textbf{Constraints:} $1 \le t \le 100$ test cases, $1 \le n, m \le 50$.

\begin{table}[H]
\centering\small
\begin{tabular}{@{}clll@{}}
\toprule
\textbf{\#} & \textbf{Maze} & \textbf{Official / our output} & \textbf{Why} \\
\midrule
1 & \texttt{.} & \texttt{Yes} / \texttt{Yes} & nobody there \\
2 & \texttt{G.} & \texttt{Yes} / \texttt{Yes} & G is next to the exit \\
3 & \texttt{\#B / G.} & \texttt{No} / \texttt{No} & walling B in walls the exit, G is stuck \\
4 & \texttt{G.\# / B\#.} & \texttt{No} / \texttt{No} & G stands next to B \\
5 & \texttt{\#B. / \#.. / GG.} & \texttt{Yes} / \texttt{Yes} & wall the two cells next to B \\
6 & \texttt{\#B / B.} & \texttt{Yes} / \texttt{Yes} & no G: wall the exit itself \\
\bottomrule
\end{tabular}
\end{table}

\begin{diagrambox}
\begin{tikzpicture}[x=6.5mm, y=6.5mm]
    \begin{scope}
        \gptworow{0}{w/, B/}\gptworow{1}{G/, e/}
        \node[font=\scriptsize\bfseries] at (0.5,1.1) {sample 3};
        \node[note] at (0.5,-2.1) {No};
    \end{scope}
    \begin{scope}[xshift=2.3cm]
        \gptworow{0}{G/, e/, w/}\gptworow{1}{B/, w/, e/}
        \node[font=\scriptsize\bfseries] at (1,1.1) {sample 4};
        \node[note] at (1,-2.1) {No};
    \end{scope}
    \begin{scope}[xshift=5.2cm]
        \gptworow{0}{w/, B/, e/}\gptworow{1}{w/, e/, e/}\gptworow{2}{G/, G/, e/}
        \node[font=\scriptsize\bfseries] at (1,1.1) {sample 5};
        \node[note] at (1,-2.8) {Yes};
    \end{scope}
    \begin{scope}[xshift=8.1cm]
        \gptworow{0}{w/, B/, x/}\gptworow{1}{w/, x/, o/}\gptworow{2}{G/, G/, o/}
        \node[font=\scriptsize\bfseries] at (1,1.1) {5: our walls};
        \node[note] at (1,-2.8) {G, G reached};
    \end{scope}
    \begin{scope}[xshift=11.0cm]
        \gptworow{0}{w/, B/}\gptworow{1}{B/, x/}
        \node[font=\scriptsize\bfseries] at (0.5,1.1) {sample 6};
        \node[note] at (0.5,-2.1) {Yes (exit walled)};
    \end{scope}
\end{tikzpicture}
\figcap{Some samples (exit $=$ bottom-right cell; grey $=$ wall, red frame $=$ a wall we add, blue $=$ reached from the exit). In sample 5 the cells right of B and below B are walled; the search from the exit then reaches both G.}
\end{diagrambox}

\subsubsection*{2. The graph model}

\begin{modelbox}[Graph Model --- Solve The Maze]
\renewcommand{\arraystretch}{1.35}
\begin{tabularx}{\linewidth}{@{}>{\bfseries}l X@{}}
Vertices & the non-wall cells, up to $50 \cdot 50 = 2500$ per test \\
Edges & undirected, between side-adjacent non-wall cells; implicit (computed from the grid) \\
Weights & none \\
Question & choose a set $W$ of empty cells to delete so that in the remaining graph every G is in the component of the exit and no B is --- equivalently: one search \textbf{from the exit} marks all G and no B \\
Size & $2500$ vertices, $\le 5000$ edges per test, 100 tests: any linear algorithm is instant; trying sets $W$ is not ($2^{2500}$) \\
\end{tabularx}
\end{modelbox}

\begin{diagrambox}
\begin{tikzpicture}[x=9mm, y=9mm]
    \begin{scope}
        \gptworow{0}{w/, B/, e/}\gptworow{1}{w/, e/, e/}\gptworow{2}{G/, G/, e/}
        \node[font=\scriptsize\bfseries] at (1,0.95) {sample 5 as a grid};
    \end{scope}
    \begin{scope}[xshift=4.2cm]
        \foreach \r/\c/\l in {0/1/B, 0/2/{}, 1/1/{}, 1/2/{}, 2/0/G, 2/1/G, 2/2/{}}
            \node[gv, minimum size=6mm, font=\scriptsize\bfseries] (c\r\c) at (\c,-\r) {\l};
        \foreach \a/\b in {01/02, 01/11, 02/12, 11/12, 11/21, 12/22, 20/21, 21/22} \draw[ge] (c\a) -- (c\b);
        \node[gvdst, minimum size=6mm, font=\scriptsize\bfseries] at (2,-2) {ex};
        \node[font=\scriptsize\bfseries] at (1,0.95) {the cell graph};
    \end{scope}
    \node[note, text width=5.2cm, align=left, anchor=west] at (8.3,-1) {7 open cells, 8 edges. B is connected to the exit through $(1,3)$ and through $(2,2)$: both must be deleted --- they are exactly the empty neighbours of B. What remains of the exit's component is $\{(2,3), (3,3), (3,2), (3,1)\}$, containing both G.};
\end{tikzpicture}
\figcap{The model: cells are vertices, sides are edges, walls are missing vertices. Adding a wall deletes a vertex.}
\end{diagrambox}

\subsubsection*{3. Key observations}

\begin{itemize}[leftmargin=*]
    \item \textbf{Reverse the roles.} ``Every G can reach the exit'' is one question per G; ``the exit reaches every G'' is one search. The graph is undirected, so nothing has to be reversed except the starting point.
    \item \textbf{A G next to a B is fatal.} If G escapes, B can step onto G's cell (it cannot be walled) and follow G's route. Answer \texttt{No}.
    \item \textbf{A G route that passes next to a B is useless.} If a route from G to the exit visits a cell $c$ adjacent to B, then B can step to $c$ and follow the rest of the route out. So in \emph{every} valid solution, every escaping route avoids all neighbours of all B.
    \item \textbf{Therefore: wall all empty neighbours of every B.} It costs nothing (no valid route uses them) and it traps every B (its four sides are now walls, borders, or other B --- never G, see above). After that, no more walls are needed: just check with one search from the exit.
    \item \textbf{No G at all:} answer \texttt{Yes} --- wall the exit (allowed) or simply rely on the B being walled in. Do not run into the trap of answering \texttt{No} because the exit itself became a wall.
\end{itemize}

\begin{diagrambox}
\begin{tikzpicture}[x=6.5mm, y=6.5mm]
    \begin{scope}
        \gptworow{0}{G/, e/, B/, e/}
        \node[font=\scriptsize\bfseries] at (1.5,1.0) {\texttt{G.B.}: the exit is next to B};
        \gptworow{2}{G/, e/, B/, x/}
        \node[note, text width=3.6cm, anchor=north] at (1.5,-2.7) {walling B in walls the exit: \texttt{No}};
    \end{scope}
    \begin{scope}[xshift=4.6cm]
        \gptworow{0}{G/, e/, e/, w/}\gptworow{1}{e/, w/, e/, e/}\gptworow{2}{e/, B/, e/, e/}
        \node[font=\scriptsize\bfseries] at (1.5,1.0) {two routes of G};
        \draw[gpath, -{Stealth[length=2mm]}] (0.15,0.15) -- (2.15,0.15) -- (2.15,-0.85) -- (3.15,-0.85) -- (3.15,-1.8);
        \draw[gback] (-0.15,-0.15) -- (-0.15,-2.15) -- (2.9,-2.15);
        \node[note, text width=3.8cm, anchor=north] at (1.5,-2.7) {the bottom route (red) passes next to B, so B could follow it out: useless. The top route (orange) avoids B's neighbours};
    \end{scope}
    \begin{scope}[xshift=9.6cm]
        \gptworow{0}{G/, e/, e/, w/}\gptworow{1}{e/, w/, e/, e/}\gptworow{2}{x/, B/, x/, e/}
        \node[font=\scriptsize\bfseries] at (1.5,1.0) {wall in B};
        \node[note, text width=3.6cm, anchor=north] at (1.5,-2.7) {G still escapes along the top: \texttt{Yes}};
    \end{scope}
\end{tikzpicture}
\figcap{Why walling the neighbours of B is never a mistake: a cell next to B is useless for G, because B could use it too.}
\end{diagrambox}

\subsubsection*{4. Approaches}

\begin{bruteforcebox}
\textbf{Idea: try every set of walls.} For each subset of the empty cells, wall it, run one search from the exit (unless it is walled) and check the condition. Only for tiny grids ($\le 15$ empty cells), but a perfect reference: it uses no insight at all.

\begin{lstlisting}[style=cpp]
#include <bits/stdc++.h>
using namespace std;

int n, m;
const int dr[] = {1, -1, 0, 0}, dc[] = {0, 0, 1, -1};

bool ok(const vector<string>& h) {               // search from the exit, check G and B
    vector<vector<char>> seen(n, vector<char>(m, 0));
    if (h[n - 1][m - 1] != '#') {
        vector<pair<int, int>> st(1, make_pair(n - 1, m - 1));
        seen[n - 1][m - 1] = 1;
        while (!st.empty()) {
            int r = st.back().first, c = st.back().second; st.pop_back();
            for (int d = 0; d < 4; d++) {
                int a = r + dr[d], b = c + dc[d];
                if (a < 0 || b < 0 || a >= n || b >= m) continue;
                if (seen[a][b] || h[a][b] == '#') continue;
                seen[a][b] = 1; st.push_back(make_pair(a, b));
            }
        }
    }
    for (int r = 0; r < n; r++)
        for (int c = 0; c < m; c++) {
            if (h[r][c] == 'G' && !seen[r][c]) return false;
            if (h[r][c] == 'B' && seen[r][c]) return false;
        }
    return true;
}

int main() {
    int t;
    cin >> t;
    while (t--) {
        cin >> n >> m;
        vector<string> g(n);
        for (auto& s : g) cin >> s;
        vector<pair<int, int>> e;                // the empty cells
        for (int r = 0; r < n; r++)
            for (int c = 0; c < m; c++)
                if (g[r][c] == '.') e.push_back(make_pair(r, c));
        bool any = false;
        for (long long mask = 0; mask < (1LL << e.size()) && !any; mask++) {
            vector<string> h = g;
            for (size_t i = 0; i < e.size(); i++)
                if (mask >> i & 1) h[e[i].first][e[i].second] = '#';
            any = ok(h);
        }
        cout << (any ? "Yes" : "No") << "\n";
    }
}
\end{lstlisting}

\begin{complexitybox}
\textbf{Time:} $2^{E} \cdot \bigO{nm}$ for $E$ empty cells. Fine for the samples ($E \le 4$); a $50 \times 50$ maze can have $E = 2500$. \textbf{Space:} \bigO{nm}. \textbf{Verdict:} \tle{} beyond $E \approx 20$; our reference for random tests.
\end{complexitybox}
\end{bruteforcebox}

\begin{algorithmbox}[Approach 2: wall the bad, one search from the exit]
\textbf{Idea.} (1) For every B, look at its four neighbours: a G means \texttt{No}; an empty cell becomes a wall. (2) No G: \texttt{Yes}. (3) The exit is now a wall: \texttt{No}. (4) Search from the exit; \texttt{Yes} iff every G and no B is reached.

\begin{lstlisting}[style=cpp]
#include <bits/stdc++.h>
using namespace std;

// Solve The Maze (CF 1365D): wall in every B, then one DFS from the exit
int n, m;
vector<string> g;
const int dr[] = {1, -1, 0, 0}, dc[] = {0, 0, 1, -1};

bool inside(int r, int c) { return r >= 0 && c >= 0 && r < n && c < m; }

bool solve() {
    int goods = 0;
    for (int r = 0; r < n; r++)
        for (int c = 0; c < m; c++) {
            if (g[r][c] == 'G') goods++;
            if (g[r][c] != 'B') continue;
            for (int d = 0; d < 4; d++) {
                int nr = r + dr[d], nc = c + dc[d];
                if (!inside(nr, nc)) continue;
                if (g[nr][nc] == 'G') return false;        // B would follow G out
                if (g[nr][nc] == '.') g[nr][nc] = '#';     // block B in
            }
        }
    if (goods == 0) return true;                           // nobody must escape
    if (g[n - 1][m - 1] == '#') return false;              // the exit itself was walled
    vector<vector<char>> seen(n, vector<char>(m, 0));
    vector<pair<int, int>> st(1, make_pair(n - 1, m - 1));
    seen[n - 1][m - 1] = 1;
    int reachedG = 0, reachedB = 0;
    while (!st.empty()) {                                  // search from the exit
        int r = st.back().first, c = st.back().second; st.pop_back();
        if (g[r][c] == 'G') reachedG++;
        if (g[r][c] == 'B') reachedB++;
        for (int d = 0; d < 4; d++) {
            int nr = r + dr[d], nc = c + dc[d];
            if (!inside(nr, nc) || seen[nr][nc] || g[nr][nc] == '#') continue;
            seen[nr][nc] = 1;
            st.push_back(make_pair(nr, nc));
        }
    }
    return reachedG == goods && reachedB == 0;
}

int main() {
    int t;
    cin >> t;
    while (t--) {
        cin >> n >> m;
        g.assign(n, "");
        for (auto& s : g) cin >> s;
        cout << (solve() ? "Yes" : "No") << "\n";
    }
}
\end{lstlisting}

\begin{complexitybox}
\textbf{Time:} \bigO{nm} per test: the B scan looks at $4nm$ neighbours, the search visits each cell once. 100 tests of $50 \times 50$ cells (all G) ran in 0.2\,s including I/O. \textbf{Space:} \bigO{nm}. \textbf{Verdict:} fast; the stack holds at most $nm = 2500$ cells, and a recursive DFS (depth $\le 2500$) would be safe too.
\end{complexitybox}
\end{algorithmbox}

\begin{invariantbox}[Why it is correct]
Let $W_0$ be the set of empty neighbours of all B, and suppose no G is adjacent to a B.

\emph{If any valid wall set $W$ exists, then $W_0$ is valid.} With $W_0$ every B is trapped: its neighbours are walls, borders or B cells, so its component contains only B cells and never the (empty) exit. Now take a good person $g$. In the solution $W$ it has a route $P$ to the exit. No cell of $P$ is adjacent to a B --- otherwise that B steps onto $P$ and follows it out, contradicting the validity of $W$. So no cell of $P$ is in $W_0$, and $P$ is still open after walling $W_0$: $g$ escapes. Hence $W_0$ is valid.

\emph{If $W_0$ is not valid, nothing is.} By the contrapositive of the above. And if some G is adjacent to a B, no $W$ is valid (B follows G). If there is no G, $W_0$ is valid (every B is trapped) --- even when it walls the exit.

\emph{The search decides $W_0$.} With $W_0$ in place, the component of the exit is exactly the set marked by the search (Chapter~1), so ``every G marked, no B marked'' is the definition of valid. \hfill$\blacksquare$
\end{invariantbox}

\begin{alternativebox}[Approach 3: components with a DSU]
After walling $W_0$, the question is only ``which component is each person in?'', so any connectivity method works. A DSU over the open cells unites every pair of adjacent open cells; then every G must share the exit's root and no B may.

\begin{lstlisting}[style=cpp]
#include <bits/stdc++.h>
using namespace std;

int n, m;
vector<string> g;
vector<int> par;
int fnd(int x) { while (par[x] != x) x = par[x] = par[par[x]]; return x; }

bool solve() {
    const int dr[] = {1, -1, 0, 0}, dc[] = {0, 0, 1, -1};
    int goods = 0;
    for (int r = 0; r < n; r++)
        for (int c = 0; c < m; c++) {
            if (g[r][c] == 'G') goods++;
            if (g[r][c] != 'B') continue;
            for (int d = 0; d < 4; d++) {
                int a = r + dr[d], b = c + dc[d];
                if (a < 0 || b < 0 || a >= n || b >= m) continue;
                if (g[a][b] == 'G') return false;
                if (g[a][b] == '.') g[a][b] = '#';
            }
        }
    if (goods == 0) return true;
    if (g[n - 1][m - 1] == '#') return false;
    par.resize(n * m);
    iota(par.begin(), par.end(), 0);
    for (int r = 0; r < n; r++)                            // unite right and down
        for (int c = 0; c < m; c++) {
            if (g[r][c] == '#') continue;
            int x = r * m + c;
            if (r + 1 < n && g[r + 1][c] != '#') par[fnd(x)] = fnd(x + m);
            if (c + 1 < m && g[r][c + 1] != '#') par[fnd(x)] = fnd(x + 1);
        }
    int ex = fnd(n * m - 1);
    for (int r = 0; r < n; r++)
        for (int c = 0; c < m; c++) {
            if (g[r][c] == 'G' && fnd(r * m + c) != ex) return false;
            if (g[r][c] == 'B' && fnd(r * m + c) == ex) return false;
        }
    return true;
}

int main() {
    int t;
    cin >> t;
    while (t--) {
        cin >> n >> m;
        g.assign(n, "");
        for (auto& s : g) cin >> s;
        cout << (solve() ? "Yes" : "No") << "\n";
    }
}
\end{lstlisting}

\begin{center}
\begin{tikzpicture}[x=6.5mm, y=6.5mm]
    \gptworow{0}{Gn/, on/, x/, on/, on/}
    \gptworow{1}{on/, w/, Bn/, x/, on/}
    \gptworow{2}{on/, on/, w/, on/, Gn/}
    \gptworow{3}{on/, on/, on/, on/, on/}
    \draw[gtree] (0,-3) -- (4,-3); \draw[gtree] (0,0) -- (0,-3); \draw[gtree] (4,0) -- (4,-3);
    \draw[gtree] (0,0) -- (1,0); \draw[gtree] (3,0) -- (4,0); \draw[gtree] (0,-2) -- (1,-2);
    \draw[gtree] (1,-2) -- (1,-3); \draw[gtree] (3,-2) -- (4,-2); \draw[gtree] (3,-2) -- (3,-3);
    \node[note, text width=6.4cm, align=left, anchor=west] at (5.4,-1.5) {The trace maze of Section~5 after walling: the DSU joins all 15 open cells (blue lines: the adjacencies it unites) into the exit's set, both G included. B is alone in its own set.};
\end{tikzpicture}
\end{center}

\begin{complexitybox}
\textbf{Time:} \bigO{nm\,\alpha(nm)} per test. \textbf{Space:} \bigO{nm}. \textbf{Verdict:} same answers as Approach~2 on all tests below; more code for no gain here, but the DSU form is the one to use when walls are \emph{removed} one by one (Chapter~6).
\end{complexitybox}
\end{alternativebox}

\begin{tricksbox}[Approach 4: Implementation Refinements]
\begin{itemize}[leftmargin=*]
    \item \textbf{Check G before walling the exit.} The order ``no G $\Rightarrow$ \texttt{Yes}'' before ``exit walled $\Rightarrow$ \texttt{No}'' matters: sample 6 has a walled exit and the answer \texttt{Yes}.
    \item \textbf{Counting instead of a second scan.} Count G while walling; the search counts the G it reaches; compare the two numbers.
    \item \textbf{Modifying the grid in place} is fine (each test reads a new grid). If the grid must survive, copy it first.
    \item \textbf{The ``no B reached'' test is a safety net.} After the walling a B can only be connected to other B, never to the exit, so \texttt{reachedB} is always 0 in Approach~2 --- but the check costs nothing and protects against a bug in the walling loop.
\end{itemize}
\end{tricksbox}

\subsubsection*{5. Trace}

\begin{tracebox}[{A $4 \times 5$ maze: \texttt{G.... / .\#B.. / ..\#.G / .....}}]
Step 1 walls the empty neighbours of B $= (2,3)$: the cells $(1,3)$ and $(2,4)$ (the cells $(2,2)$ and $(3,3)$ are walls already). Two G, no G next to B, the exit is open. Step 2 is the search from the exit, neighbours tried in the order down, up, right, left. The numbers are the order in which cells are popped:
\begin{center}
\begin{tikzpicture}[x=6.5mm, y=6.5mm]
    \begin{scope}
        \gptworow{0}{G/, e/, e/, e/, e/}\gptworow{1}{e/, w/, B/, e/, e/}\gptworow{2}{e/, e/, w/, e/, G/}\gptworow{3}{e/, e/, e/, e/, e/}
        \node[font=\scriptsize\bfseries] at (2,1.0) {input};
    \end{scope}
    \begin{scope}[xshift=4.2cm]
        \gptworow{0}{G/, e/, x/, e/, e/}\gptworow{1}{e/, w/, B/, x/, e/}\gptworow{2}{e/, e/, w/, e/, G/}\gptworow{3}{e/, e/, e/, e/, e/}
        \node[font=\scriptsize\bfseries] at (2,1.0) {1. wall in B};
    \end{scope}
    \begin{scope}[xshift=8.4cm]
        \gptworow{0}{Gn/8, on/9, x/, on/15, on/14}
        \gptworow{1}{on/7, w/, Bn/B, x/, on/13}
        \gptworow{2}{on/6, on/10, w/, on/11, Gn/12}
        \gptworow{3}{on/5, on/4, on/3, on/2, on/1}
        \node[font=\scriptsize\bfseries] at (2,1.0) {2. search from the exit};
    \end{scope}
\end{tikzpicture}
\end{center}
\begin{center}\small
\begin{tabular}{@{}cll@{}}
\toprule
\textbf{pop} & \textbf{cell} & \textbf{stack after (top on the right)} \\
\midrule
1 & $(4,5)$ exit & $(3,5)\ (4,4)$ \\
2 & $(4,4)$ & $(3,5)\ (3,4)\ (4,3)$ \\
3 & $(4,3)$ & $(3,5)\ (3,4)\ (4,2)$\quad ($(3,3)$ is a wall) \\
4--5 & $(4,2)$, $(4,1)$ & $(3,5)\ (3,4)\ (3,2)\ (3,1)$ \\
6--7 & $(3,1)$, $(2,1)$ & $(3,5)\ (3,4)\ (3,2)\ (1,1)$ \\
8 & $(1,1)$ \textbf{G} & $(3,5)\ (3,4)\ (3,2)\ (1,2)$ \\
9--11 & $(1,2)$, $(3,2)$, $(3,4)$ & $(3,5)$\quad ($(1,3)$ is our wall) \\
12 & $(3,5)$ \textbf{G} & $(2,5)$ \\
13--15 & $(2,5)$, $(1,5)$, $(1,4)$ & empty \\
\bottomrule
\end{tabular}
\end{center}
Both G are reached (pops 8 and 12), B is not: \texttt{Yes}. With one more B in the corner, \texttt{G.... / .\#B.. / ..\#.G / B....}, the cells $(3,1)$ and $(4,2)$ are walled as well, the G in $(1,1)$ is cut off and the program prints \texttt{No}:
\begin{center}
\begin{tikzpicture}[x=6.5mm, y=6.5mm]
    \gptworow{0}{G/, e/, x/, on/8, on/7}
    \gptworow{1}{e/, w/, Bn/B, x/, on/6}
    \gptworow{2}{x/, e/, w/, on/4, Gn/5}
    \gptworow{3}{Bn/B, x/, on/3, on/2, on/1}
    \node[note, text width=6.8cm, align=left, anchor=west] at (5.4,-1.5) {The search from the exit pops only 8 cells. The region $\{(1,1), (1,2), (2,1)\}$ with the G is closed by our walls $(1,3)$ and $(3,1)$ --- and both are neighbours of a B, so no solution could keep them open. The cell $(3,2)$ is isolated.};
\end{tikzpicture}
\end{center}
\end{tracebox}

\subsubsection*{6. Test suite}

\begin{testbox}
\centering\small
\begin{tabular}{@{}p{4.6cm}lp{6.6cm}@{}}
\toprule
\textbf{Maze} & \textbf{Ours} & \textbf{Why} \\
\midrule
\texttt{.} & \texttt{Yes} & $1 \times 1$, nobody \\
\texttt{GB.} & \texttt{No} & G next to B \\
\texttt{B. / ..} & \texttt{Yes} & no G; B walled in \\
\texttt{G.B / ...} & \texttt{No} & B next to the exit: the exit is walled, G exists \\
\texttt{G\#. / .\#.} & \texttt{No} & G separated by walls, no B at all \\
\texttt{G.B.} & \texttt{No} & B stands between G and the exit \\
the $4 \times 5$ trace maze & \texttt{Yes} & G escapes around B \\
same with \texttt{B} at $(4,1)$ & \texttt{No} & the second B cuts G off \\
100 tests $50 \times 50$, all G (half with one B) & 50 \texttt{Yes}, 50 \texttt{No} & 0.2\,s \\
\bottomrule
\end{tabular}
\end{testbox}
All outputs come from running Approach~2. Approaches 2 and 3 were also compared with the brute force on $40\,000$ random grids of size up to $4 \times 4$ (cells \texttt{.}/\texttt{\#}/\texttt{G}/\texttt{B} with probabilities $0.5/0.2/0.2/0.1$, the exit forced empty): no difference.

\subsubsection*{7. Common mistakes}

\begin{warningbox}[Mistake 1 --- answering \texttt{No} when the exit gets walled but there is no G]
Sample 6 (\texttt{\#B / B.}): the exit is a neighbour of B and becomes a wall, but nobody needs to escape. Test ``no G'' first.
\end{warningbox}

\begin{warningbox}[Mistake 2 --- searching from every G]
Correct, but $(nm)^2 = 6.25 \cdot 10^6$ cell visits per test and up to $6 \cdot 10^8$ over 100 tests. The one search from the exit is $nm$.
\end{warningbox}

\begin{warningbox}[Mistake 3 --- walling only the cells that lie between B and the exit]
``Block the shortest path of B'' or ``block where B's search meets G's search'' leaves other routes open or wastes cells G needs. The neighbours of B are the only cells whose walling is provably free.
\end{warningbox}

\begin{warningbox}[Mistake 4 --- treating G and B cells as obstacles]
People can be walked through. A search that stops at G cells misses a G that is only reachable through another G (sample 5: the left G is reached through the right one).
\end{warningbox}

\subsubsection*{8. Takeaways}

\begin{itemize}[leftmargin=*]
    \item ``Everyone good must reach the exit'' $=$ one search from the exit.
    \item Make the bad starts harmless with the change every solution needs anyway; then only check.
    \item The proof idea --- ``a cell next to B is useless for G'' --- is the whole problem.
    \item Edge cases live at the exit: it can be walled, and with no G that is fine.
\end{itemize}

\begin{variantsbox}
\variant{Several exits}{\varnone}
  {every border cell marked \texttt{E} is an exit.}
  {same graph; the target is a set.}
  {wall in B, then one search from \emph{all} open exits at once (template 1 of the pattern).}
\variant{Minimum number of walls}{\varnone}
  {as few walls as possible.}
  {now walling all of B's neighbours may be wasteful: a vertex cut between B and the exit.}
  {minimum vertex cut, i.e.\ maximum flow on a split-vertex graph (Chapter~16); the search alone no longer suffices.}
\variant{Directed corridors}{\varnone}
  {some cells are one-way (arrows).}
  {directed cell graph.}
  {the escape check becomes a search from the exit in the \emph{reversed} graph; the ``B follows G'' argument still holds for cells B can reach.}
\variant{Enclosed regions}{\varin{https://leetcode.com/problems/surrounded-regions/}}
  {no people; find the cells that cannot reach the border.}
  {the targets are the whole border.}
  {one search from all border cells (Exercise E5.1).}
\end{variantsbox}
